\clearpage\pagestyle{otMain}\fontsize{11}{15}\selectfont\corrigesChapitre{10}{Statistiques et représentations}
\begin{blocLivre}{cor-10-1}\begin{corrige}{1}{Compter}\demarche{Compter les observations, pas seulement les valeurs différentes.}
On compte $\res6$ observations ; 4 apparaît $\res2$ fois. Les répétitions comptent comme des observations distinctes.
\end{corrige}\end{blocLivre}
\begin{blocLivre}{cor-10-2}\begin{corrige}{2}{Une fréquence}\demarche{Identifier le total de référence.}
La fréquence est $12/30=\res{0{,}4}$, soit $\res{40\,\%}$. Le dénominateur est l’effectif de tout le groupe.
\end{corrige}\end{blocLivre}
\begin{blocLivre}{cor-10-3}\begin{corrige}{3}{Une moyenne simple}\demarche{Additionner puis diviser par l’effectif.}
La somme vaut 40. La moyenne vaut $40/4=\res{10}$. Elle n’est pas nécessairement l’une des valeurs, même si c’est le cas ici.
\end{corrige}\end{blocLivre}
\begin{blocLivre}{cor-10-4}\begin{corrige}{4}{Ordonner avant la médiane}\demarche{Ranger avant de choisir le rang central.}
La liste ordonnée est $4;7;9;11;13$. La troisième des cinq valeurs est $\res9$. Deux valeurs sont de chaque côté.
\end{corrige}\end{blocLivre}
\begin{blocLivre}{cor-10-5}\begin{corrige}{5}{Mesurer l’étendue}\demarche{Soustraire le minimum, même s’il est négatif.}
Le maximum est 8 et le minimum $-3$. L’étendue vaut $8-(-3)=\res{11\,^\circ\mathrm C}$. Elle est positive.
\end{corrige}\end{blocLivre}
\begin{blocLivre}{cor-10-6}\begin{corrige}{6}{Un tableau complet}\demarche{Utiliser les effectifs cumulés pour les rangs centraux.}
L’effectif est $2+3+6+5+4=20$. La somme pondérée vaut $12+24+60+60+64=220$, d’où une moyenne de $\res{11}$. Les rangs 10 et 11 portent tous deux 10 : médiane $\res{10}$. L’étendue vaut $16-6=\res{10}$.
\end{corrige}\end{blocLivre}
\begin{blocLivre}{cor-10-7}\begin{corrige}{7}{Une médiane paire}\demarche{Distinguer un rang d’une valeur.}
Les deux valeurs centrales sont 6 et 8. La médiane retenue est $(6+8)/2=\res7$. La moyenne vaut $48/6=\res8$. La valeur 20 tire la moyenne vers le haut.
\end{corrige}\end{blocLivre}
\begin{blocLivre}{cor-10-8}\begin{corrige}{8}{Un disque statistique}\demarche{Multiplier les fréquences par un tour complet.}
Chaque réponse représente $360/40=9^\circ$. Les angles sont $\res{90^\circ}$, $\res{162^\circ}$ et $\res{108^\circ}$. Leur somme est 360 degrés.
\label{fig-cor-10-8}\figureGeo[corrige]{stat-disque}
\end{corrige}\end{blocLivre}
\begin{blocLivre}{cor-10-9}\begin{corrige}{9}{Un histogramme}\demarche{Cumuler les effectifs et tenir compte du regroupement.}
Les effectifs sont 4, 8, 6 et 2 : total $\res{20}$. Les rangs 10 et 11 sont dans la deuxième classe, donc la médiane est dans $\res{[10;20[}$. La figure ne permet pas sa valeur exacte.
\label{fig-cor-10-9}\figureGeo[corrige]{stat-histogramme}
\end{corrige}\end{blocLivre}
\begin{blocLivre}{cor-10-10}\begin{corrige}{10}{Une formule de tableur}\demarche{Pondérer avant de diviser.}
On peut saisir =A2*B2+A3*B3+A4*B4, puis diviser par SOMME(B2:B4). La somme pondérée vaut $12+35+24=71$ ; la moyenne vaut $\res{7{,}1}$. MOYENNE(A2:A4) ignorerait les effectifs.
\end{corrige}\end{blocLivre}
\begin{blocLivre}{cor-10-11}\begin{corrige}{11}{Réunir deux groupes}\demarche{Reconstituer les sommes de chaque groupe.}
Les sommes valent $12\times13=156$ et $18\times11=198$. La moyenne commune vaut $(156+198)/30=\res{11{,}8}$. Faire $(13+11)/2$ donnerait un poids égal à des groupes de tailles différentes.
\end{corrige}\end{blocLivre}
\begin{blocLivre}{cor-10-12}\begin{corrige}{12}{Une note manquante}\demarche{Traduire la moyenne en somme.}
La somme des quatre notes vaut $4\times12{,}5=50$. Les trois connues totalisent 34 ; $\res{x=16}$. Le contrôle donne bien $50/4=12{,}5$.
\end{corrige}\end{blocLivre}
\begin{blocLivre}{cor-10-13}\begin{corrige}{13}{Même moyenne, mêmes résultats ?}\demarche{Comparer un indicateur de position et un de dispersion.}
Les deux sommes valent 40 : les moyennes sont 10. Les deux médianes sont 10. Les étendues sont $\res2$ et $\res{16}$. B est beaucoup plus dispersée selon l’étendue ; les seuls indicateurs centraux masquent cet écart.
\end{corrige}\end{blocLivre}
\begin{blocLivre}{cor-10-14}\begin{corrige}{14}{Ajouter une observation}\demarche{Séparer ce que détermine la somme de ce qui demande l’ordre.}
L’ancienne somme est 126 ; la nouvelle vaut 150 pour 10 valeurs. La moyenne est $\res{15}$. La médiane reste inconnue : la somme ne renseigne pas sur les valeurs centrales. Par exemple neuf 14 et la série huit 13 plus 22 ont la même moyenne initiale mais des médianes différentes après l’ajout.
\end{corrige}\end{blocLivre}
\begin{blocLivre}{cor-10-15}\begin{corrige}{15}{Un graphique trompeur}\demarche{Ramener le pourcentage à la valeur initiale.}
La hausse réelle est $(84-80)/80=\res{5\,\%}$. Les hauteurs visibles sont 2 et 6, dans un rapport 3. Ce rapport visuel ne traduit pas un triplement des inscriptions. Lire l’origine et les graduations avant de comparer.
\label{fig-cor-10-15}\figureGeo[corrige]{stat-axe-tronque}
\end{corrige}\end{blocLivre}
